\documentclass[a4paper,11pt]{article}
\usepackage[top=2.2cm,bottom=2cm,left=2.2cm,right=2cm]{geometry}
\usepackage{fontspec}
\setmainfont{Arial}[Extension=.ttf,UprightFont=arial,BoldFont=arialbd,ItalicFont=ariali,BoldItalicFont=arialbi,Path=C:/Windows/Fonts/]
\usepackage{longtable,array}
\setlength{\parindent}{0pt}\setlength{\parskip}{6pt}
\pagestyle{empty}
\begin{document}

\begin{center}{\large\bfseries Carta --- Regra 26.4 do PCT (folhas de substituição dos desenhos)}\end{center}

{\small\itshape Minuta preparada pela Central de Patentes em 27/10/2026 (rodada 2, após revisão), para liberação do operador. Não enviada. Assinatura, local e data são do operador; a suficiência jurídica é do agente de PI [LEGAL].}

[local], [data]

\textbf{Ao Instituto Nacional da Propriedade Industrial (INPI) --- Organismo Receptor (RO/BR)}\\ Rua Mayrink Veiga, 9, 20º andar, Centro, Rio de Janeiro/RJ, CEP [dado pessoal removido]

\textbf{Pedido internacional:} PCT/BR2026/050576\\ \textbf{Ref. do processo do requerente:} BR102026005477\\ \textbf{Data do depósito internacional:} 04 de outubro de 2026\\ \textbf{Requerente:} ROTOLI MIGUEL, Luiz Antonio

\textbf{Assunto:} resposta à Solicitação de Correção de Irregularidades (PCT/RO/106) expedida em 09 de outubro de 2026 --- folhas de substituição dos desenhos (Regra 26.4 do PCT).

\textbf{1.} Em atenção ao formulário PCT/RO/106 expedido em 09 de outubro de 2026, que no Anexo C1 (desenhos do pedido internacional tal como depositado) assinala o item 1 (publicação internacional razoavelmente uniforme, Regras 11 e 26.3.a)i)), Desenhos «d» (tinta preta durável; traços de espessura uniforme e bem definidos) e «k» (números e letras de altura inferior a 0,32 cm), e remete ao formulário IB/313 [15] de 07 de outubro de 2026, Anexo C1, item 1, Desenhos «a», «d», «f», «h» e «k» («a»: não admitem reprodução direta; «d»: tinta preta durável, traços de espessura uniforme e bem definidos; «f»: distinguíveis na reprodução reduzida; «h»: números, letras e linhas de referência com clareza e simplicidade; «k»: números e letras com altura mínima de 0,32 cm), a requerente apresenta, dentro do prazo de dois meses (Regra 26.2), folhas de substituição dos desenhos. Conforme a Regra 26.4, a correção consta de folhas de substituição acompanhadas desta carta, na qual se chama a atenção para as diferenças entre as folhas substituídas e as folhas de substituição.

\textbf{2. Folhas substituídas e folhas de substituição.} São substituídas todas as 5 folhas de desenhos depositadas (1/5 a 5/5, Figuras 1 a 9). Entram 6 folhas de substituição (1/6 a 6/6), em A4, com margens de 2,5 cm (superior), 2,5 cm (esquerda), 1,5 cm (direita) e 1,0 cm (inferior), sem molduras, numeradas em algarismos arábicos no alto da folha (n/N, com N = 6), com as figuras numeradas como no depositado (Figuras 1 a 9). As folhas 1/6 (Figura 1) e 6/6 (Figura 9) estão em posição de paisagem, com o topo da figura voltado para o lado esquerdo da folha (Regra 11.13(j)).

{\small\begin{longtable}{|r|c|c|p{9.4cm}|}\hline \textbf{Fig.} & \textbf{Folha substituída} & \textbf{Folha de substituição} & \textbf{Diferenças em relação à folha substituída}\\\hline\endhead

1 & 1/5 & 1/6 & folha em paisagem (topo à esquerda, R.11.13(j)); corpo 15 pt; fundos coloridos dos domínios PALAVRA, NOME e VERBO trocados por contorno pontilhado; caixas e setas tracejadas reencaminhadas para não cruzar rótulos; as pontas de seta entre [1] e [3] e entre [4] e [5] aparecem no sentido visível no depositado ([3] para [1] e [5] para [4]; no depositado a sobreposição das caixas as deixava assim); o rótulo da caixa [1] reproduz os caracteres visíveis no depositado («VACUUM/GRAV/PHO»)\\\hline

2 & 1/5 & 2/6 & corpo 15 pt; círculos e equação central redimensionados; «Comprimir» e «Filtrar» afastados das setas, sem fundo; fundos coloridos retirados\\\hline

3 & 2/5 & 2/6 & mesma curva, mesmas fronteiras e mesmos rótulos; as três primeiras letras de «VACUUM» e parte de «PHOTON 16b», cortadas pela moldura do gráfico no depositado, aparecem inteiras (palavras da descrição, [033]); rótulos de região sem fundo branco, com a grade pontilhada interrompida ao redor deles; «VACUUM 6 bits» passou da horizontal para a vertical (mesmo texto); legenda acima do quadro; letras de corpo 15 pt\\\hline

4 & 2/5 & 3/6 & corpo 15 pt; quatro estados e eixos mantidos; w$>$0, w$<$0, $\partial$w$>$0 e $\partial$w$<$0 fora das caixas (no depositado cruzavam as caixas); fundos coloridos retirados\\\hline

5 & 3/5 & 3/6 & caixas ampliadas para conter o texto; «Próxima geração» em duas linhas abaixo do trecho horizontal da seta de retorno, sem fundo branco (a borda da caixa CARGA fica contínua); cores retiradas\\\hline

6 & 3/5 & 4/6 & colunas com hachuras distintas (GRACE oblíqua, FRONTIER oblíqua inversa, VACUUM horizontal) no lugar de tons de cinza; limiares tracejados desenhados só entre as colunas (no depositado as colunas cobriam o traço); legendas com a hachura correspondente\\\hline

7 & 4/5 & 4/6 & séries distintas por marcador e tipo de traço no lugar de cor; legenda à direita; grade pontilhada preta\\\hline

8 & 4/5 & 5/6 & colunas hachuradas no lugar de cinza; «Vácuo angular» fora do quadro, à direita, com seta oblíqua para os buckets 13 e 14 vazios, como no depositado; sem fundo branco sobre a moldura ou a grade\\\hline

9 & 5/5 & 6/6 & folha em paisagem; caixas e setas pretas; nenhuma sobreposição; siglas em maiúsculas de corpo pleno (no depositado: versaletes)\\\hline

\end{longtable}}

\textbf{3. Diferenças comuns a todas as folhas.} (a) traços e letras em preto puro sobre fundo branco, sem tons de cinza e sem cor; (b) traços de espessura uniforme e bem definidos, legíveis em reprodução reduzida a dois terços; (c) algarismos, maiúsculas e letras com haste com altura mínima de 0,32 cm (inclusive nos sinais de referência); (d) rótulos, legendas e caixas de texto posicionados sem cobrir traço: nenhum fundo branco ou elemento opaco fica sobre linha, ponto, coluna ou hachura (onde a grade pontilhada passaria por um rótulo ou legenda, a grade foi interrompida); as setas reproduzem o sentido visível no depositado; (e) folhas A4 com as margens da Regra 11.6, sem molduras; (f) disposição das figuras e dos painéis nas folhas, e número de folhas, ajustados ao tamanho das letras. Nas folhas de substituição os trechos que a moldura ou outra caixa cortava no depositado aparecem inteiros somente quando a palavra consta da descrição depositada (VACUUM e PHOTON: parágrafo [033] e legenda da Figura 3); a curva da Figura 3 foi mantida exatamente como depositada. Os cantos e margens que o Organismo Receptor usa para marcar as folhas de substituição (canto superior direito e margem inferior) ficam livres.

\textbf{4. Declaração.} As folhas de substituição corrigem somente as irregularidades formais apontadas e \textbf{não acrescentam matéria nova}: as figuras, os elementos, os sinais de referência, as palavras e os dados são os do pedido internacional tal como depositado; nenhuma figura foi acrescentada, suprimida ou renumerada. Mudaram apenas a apresentação (cor, espessura de traço, corpo das letras, disposição na folha e posição de rótulos e legendas), com as exceções declaradas no quadro do item 2.

\textbf{5.} Solicita-se que as folhas de substituição substituam integralmente as folhas de desenhos depositadas.

Anexo: FOLHAS\_SUBSTITUTIVAS\_050576.pdf (6 folhas).

Termos em que, pede deferimento.

\vspace{1.2cm}[ASSINATURA E IDENTIFICAÇÃO --- a preencher pelo operador]

\end{document}